\chapter{Calcolo deduttivo proposizionale}

Nel capitolo precedente abbiamo definito una logica come una tripla \((\mathcal{L}, \models, \vdash)\), dove \(\mathcal{L}\) è un linguaggio, \(\models\) è la relazione di soddisfacibilità che ne esprime la semantica, e \(\vdash\) è un calcolo, che ci permette di dedurre formule a partire da altre, o verificare se una formula è valida o meno.

Per quanto concerne la logica proposizionale, abbiamo definito il linguaggio \(\PL\) e fornito ben due semantiche, ma nonostante abbiamo visto una forma di calcolo tramite le tabelle di verità, questo non è sufficiente.
Infatti, questo tipo di calcolo non \qi{scala} a logiche più complesse ed espressive, poiché assume il potersi ridurre a un numero finito di casi, cosa che già per il prim'ordine non è possibile.

Vedremo in questo capitolo due calcoli, un primo detto \keyword{calcolo di deduzione naturale} più affine al ragionamento umano, e un secondo detto \keyword{calcolo di risoluzione} più affine a un'implementazione automatica, e che è alla base di molti sistemi di dimostrazione automatica.

Entrambi questi calcoli sono di tipo puramente sintattico, e non fanno riferimento alla semantica, ma le regole di trasformazione che li caratterizzano sono formulate in modo da preservare le proprietà semantiche, e quindi ci permettono di dedurre conseguenze logiche, ovvero di ottenere la stessa relazione di deducibilità che abbiamo definito tramite la semantica, ma in maniera puramente sintattica.

\begin{definition}{Calcolo inferenziale}
  Definiamo un \keyword{calcolo inferenziale} come una coppia \((\Axi, \Rules)\) dove \(\Axi \subseteq \PL\) è un insieme di formule detto \keyword{assiomi}, e \(\Rules\) è un insieme di regole di trasformazione dette \keyword{regole di inferenza}.

  Le formule in \(\Axi\) sono proposizioni considerate vere a priori, e le regole in \(\Rules\) sono regole che ci permettono a partire da un numero finito di formule, dette \keyword{premesse}, di dedurre una nuova formula, detta \keyword{conclusione}.
\end{definition}

Ovviamente, come anticipato non tutte le regole di trasformazioni sono regole d'inferenza valide, poiché vorremmo che queste regole preservassero delle proprietà semantiche.

A prescindere da ciò però, a partire da un calcolo inferenziale \((\Axi, \Rules)\) possiamo definire una relazione di deducibilità.

\begin{definition}{Deducibilità}
  Dati un calcolo inferenziale \((\Axi, \Rules)\) e un insieme di formule \(\Gamma\), diciamo che \(\Gamma\) deduce una formula \(\phi\) tramite \(\Rules\), e scriviamo \(\Gamma \vdash_{\Rules} \phi\), se è possibile dedurre come conclusione \(\phi\) a partire da un numero finito di formule in \(\Gamma\) e \(\Axi\), applicando un numero finito di regole in \(\Rules\).

  In caso non ci sia ambiguità sul calcolo inferenziale, scriveremo semplicemente \(\Gamma \vdash \phi\).
\end{definition}

Intuitivamente, l'obiettivo che vorremmo ottenere con un calcolo inferenziale, e quindi con la relazione di deducibilità, sarebbe quello di catturare, in maniera puramente sintattica, la conseguenza logica, ovvero vorremmo che:
\[\Gamma \models \phi \iff \Gamma \vdash \phi\]

Per quanto riguarda gli assiomi, sono un insieme minimale di tautologie dalle quali tramite il teorema di sostituzione varranno in diverse istansiazioni.\footnote{In un certo senso la regola di sostituzione fa parte di \(R\).}

Il primo calcolo che vedremo però non ha assiomi, ma solo regole di inferenza, e si chiama \keyword{calcolo di deduzione naturale}.

\section{Calcolo di deduzione naturale}

Avendo già anticipato l'assenza di assiomi, il calcolo di deduzione naturale esclusivamente dalle regole di inferenza, che sono formulate in modo da catturare il significato semantico degli operatori logici, e quindi ci permettono di dedurre conseguenze logiche.\footnote{
  La versione del calcolo di deduzione naturale presentato qui non è la sua versione minimale, e conterrà delle regole che non sono strettamente necessarie, ma che rendono più semplice la deduzione di formule, e quindi più affine al ragionamento umano.}

In particolare, per ogni operatore logico avremo una regola detta di \keyword{introduzione}, che ci permetterà di dedurre una formula che ha quell'operatore come connettivo principale, e una regola detta di \keyword{eliminazione}, che ci permetterà di dedurre formule che non hanno quell'operatore come connettivo principale, a partire da formule che invece lo hanno.

Nel nostro caso, consideremo i seguenti operatori logici: \(\land\), \(\lor\), \(\to\), \(\neg\) e \(\bot\), dove \(\bot\) è un operatore \(0\)-ario che rappresenta la contraddizione.

\begin{note}{}
  La notazione presentata è da intendersi come segue: sopra la linea ci sono le premesse, sotto la linea c'è la conclusione, e a destra della regola c'è il nome della regola, che è puramente indicativo, e serve solo a identificare la regola, ma non ha alcun significato formale.

  Inoltre, quando ci sono delle assunzioni che è possibile scaricare, queste sono indicate tra parentesi quadre, mentre i puntini di sospensione indicano l'applicazione di un numero finito di regole da quelle assunzioni fino alla conclusione.
\end{note}

\paragraph{Regole per \(\land\)}

Le regole di introduzione e eliminazione per \(\land\) sono le più intuitive, e sono le seguenti:

\begin{equation*}
  \infer[(\land I)]{A \land B}{A & B} \quad\quad\quad \infer[(\land E^l)]{A}{A \land B} \quad\quad\quad \infer[(\land E^r)]{B}{A \land B}
\end{equation*}

\paragraph{Regole per \(\to\)}

Per quanto riguarda \(\to\), la regola di eliminazione è abbastanza intuitiva oltre che a essere molto conosciuta col nome di \textit{Modus Ponens}, mentre la regola di introduzione è un po' più complessa.

\begin{equation*}
  \infer[(\to E)]{B}{A \to B & A} \quad\quad\quad
  \infer[(\to I)]{A \to B}{
    \begin{array}{c}
    [A] \\
    \vdots \\
    B
  \end{array}
    } 
\end{equation*}

Quello che ci dice è che, se assumendo \(A\) riusciamo a dedurre \(B\), allora possiamo dedurre \(A \to B\) senza più assumere \(A\) \keyword{scaricando l'assunzione}. Tale regola cattura il procedimento ususale di dimostrazione di un teorema che prevede un'implicazione, ovvero volendo dimostrare \(A \to B\), allora si assume \(A\) e si tenta di dedurre \(B\). In caso affermativo, si potrà concludere \(A \to B\).

\paragraph{Regole per \(\lor\)}

Le regole di introduzione per \(\lor\) sono abbastanza intuitive, mentre la regola di eliminazione è un po' più complessa.

\begin{equation*}
  \infer[(\lor I^l)]{A \lor B}{A} \quad\quad\quad \infer[(\lor I^r)]{A \lor B}{B} \quad\quad\quad
  \infer[(\lor E)]{C}{
    \begin{array}[b]{c}
    A \lor B
    \end{array} &
    \begin{array}[b]{c}
    [A] \\
    \vdots \\
    C
  \end{array} &
    \begin{array}[b]{c}
    [B] \\
    \vdots \\
    C
  \end{array}
  }
\end{equation*}

Quest'ultima regola ci dice che, se abbiamo \(A \lor B\), e se assumendo \(A\) e \(B\) \textbf{separatamente} riusciamo a dedurre \(C\) in entrambi i casi, allora possiamo dedurre \(C\) senza più assumere \(A\) e \(B\), ovvero scaricando entrambe le assunzioni nelle rispettive deduzioni.
Questa regola cattura il concetto di dimostrazione per casi, ovvero se si vuole dimostrare \(C\) a partire da \(A \lor B\), sarà mi è sufficiente dedurre \(C\) sia nel caso in cui sia vero \(A\), che nel caso in cui sia vero \(B\). In caso affermativo, si potrà concludere \(C\).
 
\paragraph{Regoler per \(\neg\)}

Le regole di introduzione e eliminazione per \(\neg\) sono le seguenti:

\begin{equation*}
  \infer[(\neg I)]{\neg A}{
    \begin{array}{c}
    [A] \\
    \vdots \\
    \bot
  \end{array}
  } \quad\quad\quad
  \infer[(\neg E)]{\bot}{\neg A & A}
\end{equation*}

La prima regola ci dice che, se assumendo \(A\) riusciamo a dedurre \(\bot\), allora possiamo dedurre \(\neg A\) senza più assumere \(A\) scaricando l'assunzione. Tale regola cattura il procedimento ususale di dimostrazione di un teorema che prevede una negazione, ovvero volendo dimostrare \(\neg A\), allora si assume \(A\) e si tenta di dedurre \(\bot\). In caso affermativo, si potrà concludere \(\neg A\).

\paragraph{Regola di riduzione all'assurdo}

Infine, vediamo due regole che potremmo definire di eliminazione di \(\bot\), di cui una generalizzazione dell'altra.

\begin{equation*}
  \infer[(\bot E)]{A}{\bot} \quad\quad\quad
  \infer[(\RAA)]{A}{
    \begin{array}{c}
    [\neg A] \\
    \vdots \\
    \bot
  \end{array}
  }
\end{equation*}

La seconda regola è esattamente la \textit{Reductio ad Absurdum}, e ci dice che, se assumendo \(\neg A\) riusciamo a dedurre \(\bot\), allora possiamo dedurre \(A\) senza più assumere \(\neg A\) scaricando l'assunzione. Tale regola cattura il procedimento ususale di dimostrazione per assurdo, in cui negando la tesi si riesce a concludere una contraddizione. L'altra regola invece, che è una generalizzazione della seconda, cattura il concetto di implicazione da premesse false.
Essenzialmente quello che esprime è che, a partire da premesse contraddittorie, è possibile dedurre qualsiasi formula, e quindi in particolare \(A\). Un analogia per avere più chiara una giustificazione intuitiva di questa regola è considerare la semantica usuale dell'implicazione, nella quale, con antecedente falso, non importa il valore di verità del conseguente, ovvero a partire da premesse false è vero tutto e il contrario di tutto.

Il processo di deduzione naturale può quindi essere visto come una costruzione incrementale: si parte da un insieme di assunzioni iniziali e, mediante passi atomici costituiti da applicazioni di regole di inferenza, si ottiene progressivamente la formula obiettivo.
Durante tale processo, alcune assunzioni vengono scaricate tramite le regole che lo consentono (ad esempio \((\to I)\), \((\neg I)\), \((\lor E)\), \((\RAA)\)), mentre le altre restano attive fino alla fine della derivazione.
Queste ultime costituiscono le assunzioni residue, e sono precisamente quelle che compaiono a sinistra del simbolo di deducibilità.
In altri termini, una derivazione conclude \(\Gamma \vdash \phi\) quando \(\phi\) è ottenuta a partire da assunzioni non scaricate esattamente pari a \(\Gamma\).

Vediamo quindi una formalizzazione più rigorosa di questo processo, e in particolare di come si costruisce un albero di deduzione.

\begin{definition}{Albero di deduzione ben formato}
  Un \keyword{albero di deduzione} è un albero finito in cui nodi sono etichettati con formule di \(\PL\).
  
  Un tale albero è detto \keyword{ben formato} se ogni nodo interno è ottenuto dai suoi figli mediante un'istanza di una regola del calcolo naturale.
  
  Le etichette delle foglie sono \keyword{occorrenze di assunzioni}.
  Se un nodo interno è ottenuto da una regola che scarica un'assunzione, allora tutte le occorrenze di quell'assunzione nel sottoalbero radicato in quel nodo sono dette \keyword{scaricate}.
\end{definition}

\begin{definition}{Deduzione nel calcolo naturale}
  Dato \(\Gamma\subseteq \PL\) e \(\phi\in\PL\), il \keyword{sequente} \(\Gamma \vdash \phi\) è detto \keyword{deducibile} se esiste un albero di deduzione ben formato con radice etichettata da \(\phi\) e foglie etichettate da formule in \(\Gamma\) o assunzioni scaricate.
\end{definition}

Andiamo ora a vedere un primo esempio di deduzione naturale, per poi procedere con esempi più complessi che ci aiuteranno a costruire l'intuizione necessaria per affrontare deduzioni più complesse, guidando la scelta delle regole da applicare.

\begin{example}[namedlabel={ex:natural_deduction}{Deduzione Naturale}]{Deduzione naturale}
  Supponiamo di voler dimostrare che \(\vdash A \to \neg \neg A\).\footnote{Indichiamo con \(\vdash A\) per indicare che \(A\) è dimostrabile senza ipotesi, ovvero che \(A\) è un teorema.}

  Chiaramente \(A \equiv \neg \neg A\), quindi ci aspettiamo di poter dimostrare questa formula.

  Iniziamo assumendo \(A\) e \(\neg A\), quindi:
  \begin{prooftree}\label{prtree:natural_deduction_example}
    \AxiomC{\({[A]}_2\)}
    \AxiomC{\({[\neg A]}_1\)}
    \BinaryInfC{\(\bot\)}
    \RuleLabel{\(\neg I_1\)}
    \UnaryInfC{\(\neg\neg A\)}
    \RuleLabel{\(\to I_2\)}
    \UnaryInfC{\(A \to \neg\neg A\)}
  \end{prooftree}

  Andiamo a etichettare le inferenze scaricate con i numeri delle regole di inferenza che le scaricano, in modo da poterle identificare facilmente.

  Possiamo vedere la rappresenzazione come albero della deduzione come segue:
  \begin{center}
    \begin{tikzpicture}
      \node (root) {\(A \to \neg\neg A\)};
      \node (n1) [above =of root] {\(\neg\neg A\)};
      \node (n2) [above =of n1] {\(\bot\)};
      \node (n3) [above left=of n2] {\([A]\)};
      \node (n4) [above right=of n2] {\([\neg A]\)};
      \draw (n1) -- (root);
      \draw (n2) -- (n1);
      \draw (n3) -- (n2);
      \draw (n4) -- (n2);
    \end{tikzpicture}
  \end{center}
\end{example}

\section{Applicazione del calcolo naturale}

In questa sezione andremo ad analizzare nel dettaglio alcuni esempi di deduzioni naturali che esploreranno tutte le regole di inferenza che abbiamo introdotto, in modo da avere un'idea più chiara di come funziona questo calcolo, e di come scegliere le regole da applicare per dedurre formule più complesse.

\subsection*{Dimostrazione di \(\vdash A \leftrightarrow \neg \neg A\).}


Sebbene nel calcolo non abbiamo considerato regole di deduzione per \(\leftrightarrow\), possiamo considerare sempre la formulazione equivalente mediante \(\land\) e \(\to\), ovvero \((A \to \neg \neg A) \land (\neg \neg A \to A)\).

Possiamo facilmente notare che per dimostrare una formula con \(\land\) come connettivo principale, è ragionevole pensare di usare \(\land I\), che ci permette di dedurre una formula con \(\land\) come connettivo principale a partire da due formule che hanno come connettivo principale i due operandi di \(\land\). Quindi come primo passo, possiamo pensare di usare \(\land I\) per ottenere \((A \to \neg \neg A) \land (\neg \neg A \to A)\) a partire da \(A \to \neg \neg A\) e \(\neg \neg A \to A\).

\begin{prooftree}
  \AxiomC{\(A \to \neg \neg A\)}
  \AxiomC{\(\neg \neg A \to A\)}
  \RuleLabel{\(\land I\)}
  \BinaryInfC{\((A \to \neg \neg A) \land (\neg \neg A \to A)\)}
\end{prooftree}

Ovviamente, \(A \to \neg \neg A\) è già stato dimostrato in \Cref{prtree:natural_deduction_example}, quindi possiamo banalmente sostituire quella foglia con l'albero di deduzione che abbiamo già costruito, e quindi ci resta da dimostrare \(\neg \neg A \to A\).

\begin{prooftree}
  \AxiomC{\Namedcref{ex:natural_deduction}}
  \UnaryInfC{\(A \to \neg \neg A\)}
  \AxiomC{\(\neg \neg A \to A\)}
  \RuleLabel{\(\land I\)}
  \BinaryInfC{\((A \to \neg \neg A) \land (\neg \neg A \to A)\)}
\end{prooftree}

Dobbiamo quindi dimostrare \(\neg \neg A \to A\). Per fare questo, possiamo andare a retroso, e pensare di usare la regola \(\to I\), che ci permette di dedurre un'implicazione a partire dal suo conseguente. Inoltre tale regola è particolarmente importante perché ci \qi{regala} un'assunzione scaricabile, ovvero ci permette di usare liberamente come assunzione il suo antecedente nel sottoalbero radicato in essa, avendo la possibilità di scaricarla alla fine della deduzione. Le assunzioni scaricabili possono essere da guida per la scelta dei passi successivi, poiché se abbiamo un'assunzione scaricabile di una formula \(A\), allora è ragionevole pensare di usare tale formula come premessa per dedurre altre formule, e quindi di usare regole di inferenza che hanno \(A\) come premessa.

\begin{prooftree}
  \AxiomC{\Namedcref{ex:natural_deduction}}
  \UnaryInfC{\(A \to \neg \neg A\)}
  \AxiomC{\(A\)}
  \RuleLabel{\(\to I_1\)}[\({[\neg\neg A]}_1\)]
  \UnaryInfC{\(\neg \neg A \to A\)}
  \RuleLabel{\(\land I\)}
  \BinaryInfC{\((A \to \neg \neg A) \land (\neg \neg A \to A)\)}
\end{prooftree}

A questo punto dobbiamo dimostrare \(A\), sapendo di poter usare liberamente \(\neg \neg A\) come assunzione. Osservando le regole che abbiamo a disposizione, nessuna ci permette direttamente di dedurre \(A\) da \(\neg \neg A\), ma \(\RAA\) ci permette di dedurre \(A\) da \(\bot\) scaricando l'assunzione di \(\neg A\).

\begin{prooftree}
  \AxiomC{\Namedcref{ex:natural_deduction}}
  \UnaryInfC{\(A \to \neg \neg A\)}
  \AxiomC{\(\bot\)}
  \RuleLabel{\(\RAA_2\)}[\({[\neg A]}_2\)]
  \UnaryInfC{\(A\)}
  \RuleLabel{\(\to I_1\)}[\({[\neg\neg A]}_1\)]
  \UnaryInfC{\(\neg \neg A \to A\)}
  \RuleLabel{\(\land I\)}
  \BinaryInfC{\((A \to \neg \neg A) \land (\neg \neg A \to A)\)}
\end{prooftree}

A questo punto il gioco è fatto, poiché ci resta da dimostrare \(\bot\), ma abbiamo come assunzioni scaricabili sia \(\neg A\) che \(A\), e quindi possiamo trivialmente usare \(\neg E\).

\begin{prooftree}
  \AxiomC{\Namedcref{ex:natural_deduction}}
  \UnaryInfC{\(A \to \neg \neg A\)}
  \AxiomC{\({[\neg\neg A]}_1\)}
  \AxiomC{\({[\neg A]}_2\)}
  \RuleLabel{\(\neg E\)}
  \BinaryInfC{\(\bot\)}
  \RuleLabel{\(\RAA_2\)}
  \UnaryInfC{\(A\)}
  \RuleLabel{\(\to I_1\)}
  \UnaryInfC{\(\neg \neg A \to A\)}
  \RuleLabel{\(\land I\)}
  \BinaryInfC{\((A \to \neg \neg A) \land (\neg \neg A \to A)\)}
\end{prooftree}

\subsection*{Dimostrazione di \((P\to Q)\to((Q\to R) \to (P\to R))\).}

Anche qui, il connettivo principale è \(\to\), che vogliamo dimostrare partendo da un insieme vuoto di premesse, quindi dovremo scaricare ogni assunzione che facciamo.
Possiamo ragionevolmente iniziare con \(\to I\) per ottenere la nostra formula a partire da \((Q\to R) \to (P\to R)\) e scaricando l'assunzione di \(P\to Q\).

\begin{prooftree}
  \AxiomC{\((Q\to R) \to (P\to R)\)}
  \RuleLabel{\(\to I_1\)}[\({[P\to Q]}_1\)]
  \UnaryInfC{\((P\to Q)\to((Q\to R) \to (P\to R))\)}
\end{prooftree}

Ci troviamo adesso a dover dimostrare un altra implicazione, e quindi possiamo pensare di usare nuovamente \(\to I\) per ottenere \((Q\to R) \to (P\to R)\) a partire da \(P\to R\) scaricando l'assunzione di \(Q\to R\).

\begin{prooftree}
  \AxiomC{\(P\to R\)}
  \RuleLabel{\(\to I_2\)}[\({[Q\to R]}_2\)]
  \UnaryInfC{\((Q\to R) \to (P\to R)\)}
  \RuleLabel{\(\to I_1\)}[\({[P\to Q]}_1\)]
  \UnaryInfC{\((P\to Q)\to((Q\to R) \to (P\to R))\)}
\end{prooftree}

Ancora una volta, ci troviamo a dover dimostrare un'altra implicazione, e quindi possiamo pensare di usare nuovamente \(\to I\) per ottenere \(P\to R\) a partire da\(R\) scaricando l'assunzione di \(P\).

\begin{prooftree}
  \AxiomC{\(R\)}
  \RuleLabel{\(\to I_3\)}[\({[P]}_3\)]
  \UnaryInfC{\(P\to R\)}
  \RuleLabel{\(\to I_2\)}[\({[Q\to R]}_2\)]
  \UnaryInfC{\((Q\to R) \to (P\to R)\)}
  \RuleLabel{\(\to I_1\)}[\({[P\to Q]}_1\)]
  \UnaryInfC{\((P\to Q)\to((Q\to R) \to (P\to R))\)}
\end{prooftree}

Adesso, facendoci guidare dalle assunzioni scaricabili, possiamo notare che usando \(\to E\) possiamo dedurre \(R\) a partire da \(Q\to R\) assumendo \(Q\).

\begin{prooftree}
  \AxiomC{\(Q\)}
  \AxiomC{\({[Q\to R]}_2\)}
  \RuleLabel{\(\to E\)}
  \BinaryInfC{\(R\)}
  \RuleLabel{\(\to I_3\)}[\({[P]}_3\)]
  \UnaryInfC{\(P\to R\)}
  \RuleLabel{\(\to I_2\)}
  \UnaryInfC{\((Q\to R) \to (P\to R)\)}
  \RuleLabel{\(\to I_1\)}[\({[P\to Q]}_1\)]
  \UnaryInfC{\((P\to Q)\to((Q\to R) \to (P\to R))\)}
\end{prooftree}

Se terminassimo ora la dimostrazione, avremmo una ottenuto una deduzione per il sequente \(Q\vdash (P\to Q)\to((Q\to R) \to (P\to R))\), poiché l'unica assunzione non scaricata è \(Q\). Usando però le assunzioni scaricabili, possiamo dedurre \(Q\) usando ancora una volta \(\to E\) a partire da \(P\to Q\) e \(P\).

\begin{prooftree}
  \AxiomC{\({[P\to Q]}_1\)}
  \AxiomC{\({[P]}_3\)}
  \RuleLabel{\(\to E\)}
  \BinaryInfC{\(Q\)}
  \AxiomC{\({[Q\to R]}_2\)}
  \RuleLabel{\(\to E\)}
  \BinaryInfC{\(R\)}
  \RuleLabel{\(\to I_3\)}
  \UnaryInfC{\(P\to R\)}
  \RuleLabel{\(\to I_2\)}
  \UnaryInfC{\((Q\to R) \to (P\to R)\)}
  \RuleLabel{\(\to I_1\)}
  \UnaryInfC{\((P\to Q)\to((Q\to R) \to (P\to R))\)}
\end{prooftree}

Solo ora, possiamo concludere la dimostrazione, poiché ogni assunzione è stata scaricata.

\subsection*{Dimostrazione di \(\neg(P\land\neg P)\).}

Il connettivo principale ora è \(\neg\), e quindi la strada più promettente è quella di usare \(\neg I\) per ottenere \(\neg(P\land\neg P)\) a partire da \(\bot\) scaricando l'assunzione di \(P\land\neg P\).

\begin{prooftree}
  \AxiomC{\(\bot\)}
  \RuleLabel{\(\neg I\)}[\({[P\land\neg P]}_1\)]
  \UnaryInfC{\(\neg(P\land\neg P)\)}
\end{prooftree}

Dobbiamo quindi dimostrare \(\bot\) avendo a disposizione \(P\land\neg P\) come assunzione scaricabile. Essendo \(\neg E\) l'unica regola che ci permette di dedurre \(\bot\), è ragionevole pensare di usarla. In particolare, per poterla applicare è necessario avere una formula \(P\) e la sua negazione.

\begin{prooftree}
  \AxiomC{\(P\)}
  \AxiomC{\(\neg P\)}
  \RuleLabel{\(\neg E\)}
  \BinaryInfC{\(\bot\)}
  \RuleLabel{\(\neg I\)}[\({[P\land\neg P]}_1\)]
  \UnaryInfC{\(\neg(P\land\neg P)\)}
\end{prooftree}

Possiamo notare però che entrambe queste assunzioni sono deducibili dalle regole di eliminazione per \(\land\) a partire da \(P\land\neg P\), che sappiamo essere un'assunzione scaricabile, e quindi possiamo usarla come premessa per dedurre sia \(P\) che \(\neg P\).

\begin{prooftree}
  \AxiomC{\({[P\land\neg P]}_1\)}
  \RuleLabel{\(\land E^l\)}
  \UnaryInfC{\(P\)}
  \AxiomC{\({[P\land\neg P]}_1\)}
  \RuleLabel{\(\land E^r\)}
  \UnaryInfC{\(\neg P\)}
  \RuleLabel{\(\neg E\)}
  \BinaryInfC{\(\bot\)}
  \RuleLabel{\(\neg I\)}[\({[P\land\neg P]}_1\)]
  \UnaryInfC{\(\neg(P\land\neg P)\)}
\end{prooftree}

Questo esempio ci mostra come un assunzione scaricabile possa essere usata più volte come premessa per dedurre formule diverse, a patto che la regola che la scarica sia radice di un sottoalbero che contiene tutte le deduzioni che la usano come premessa. In questo caso, essendo che \(\neg I\) è la radice di tutto l'albero, possiamo usare \(P\land\neg P\) liberamente, ma se tale regola fosse stata un altro nodo interno, avremmo potuto scaricarla solo nei rami che discendevano da quel nodo, e non in altri rami.

\begin{note}{}
  Le deduzioni che abbiamo visto fino a ora sono tutte deduzioni di tautologie note. Un altro modo per farci guidare nella scelta delle regole da applicare è quello di pensare alla validità semantica della formuala che vogliamo dimostrare o a eventuali equivalenze, per capire da quale assunzione \q{estrarre} l'informazione che ci manca per la dimostrazione.

  Formalemente questo processo non è corretto, poiché assume che il nostro calcolo inferenziale catturi la conseguenza logica, ovvero abbia le proprietà di correttezza e completezza, le quali, senza una dimostrazione formale, non possiamo dare per scontate.

  Per quanto riguarda il calcolo di deduzione naturale però, come dimostreremo in seguito, è possibile dimostrare che esso è corretto e completo, e quindi possiamo usare questo tipo di ragionamento per guidarci nella scelta delle regole da applicare, avendo la certezza che se una formula è valida, allora sarà dimostrabile, e viceversa.
\end{note}

\subsection*{Dimostrazione di \(P \lor \neg P\).}

In questo caso, il connettivo principale è \(\lor\), e quindi un idea ragionevole è quella di iniziare \(\lor I\) per ottenere \(P \lor \neg P\) a partire da \(P\) o \(\neg P\).
Per ora iniziamo a ottenerlo assumendo \(P\) e usando \(\lor I^l\).

\begin{prooftree}
  \AxiomC{\(P\)}
  \RuleLabel{\(\lor I^l\)}
  \UnaryInfC{\(P \lor \neg P\)}
\end{prooftree}

Non avendo alcuna assunzione scaricabile, la strada più promettente è quella di usare \(\RAA\) per ottenere \(P\) a partire da \(\bot\) scaricando l'assunzione di \(\neg P\).

\begin{prooftree}
  \AxiomC{\(\bot\)}
  \RuleLabel{\(\RAA_1\)}[\({[\neg P]}_1\)]
  \UnaryInfC{\(P\)}
  \RuleLabel{\(\lor I^l\)}
  \UnaryInfC{\(P \lor \neg P\)}
\end{prooftree}

Ora però dobbiamo dimostrare \(\bot\) avendo a disposizione \(\neg P\) come assunzione scaricabile. Come già visto, l'unica regola che ci permette di dedurre \(\bot\) è \(\neg E\), ma per poterla applicare in questo caso è necessario assumere \(P\) come premessa andando a creare un circolo vizioso.

\begin{prooftree}
  \AxiomC{\(P\)}
  \AxiomC{\({[\neg P]}_1\)}
  \RuleLabel{\(\neg E\)}
  \BinaryInfC{\(\bot\)}
  \RuleLabel{\(\RAA_1\)}
  \UnaryInfC{\(P\)}
  \RuleLabel{\(\lor I^l\)}
  \UnaryInfC{\(P \lor \neg P\)}
\end{prooftree}

Dobbiamo quindi pensare a un altro modo per dedurre \(\bot\) a partire da \(\neg P\). Come detto prima, \(\lor I\) può essere usato per dedurre \(P \lor \neg P\) a partire sia da \(P\) che da \(\neg P\). Possiamo quindi dedurre \(P \lor \neg P\) a partire da \(\neg P\) usando \(\lor I^r\), in modo da poter dedurre \(\bot\) assumendo \(\neg(P\lor\neg P)\).

\begin{prooftree}
  \AxiomC{\({[\neg P]}_1\)}
  \RuleLabel{\(\lor I^r\)}
  \UnaryInfC{\(P \lor \neg P\)}
  \AxiomC{\(\neg(P\lor\neg P)\)}
  \RuleLabel{\(\neg E\)}
  \BinaryInfC{\(\bot\)}
  \RuleLabel{\(\RAA_1\)}
  \UnaryInfC{\(P\)}
  \RuleLabel{\(\lor I^l\)}
  \UnaryInfC{\(P \lor \neg P\)}
\end{prooftree}

A questo punto abbiamo dimostrato \(\neg (P\lor \neg P)\vdash P \lor \neg P\). Per poter arrivare al teorema, ovvero \(\vdash P \lor \neg P\), dobbiamo scaricare l'assunzione di \(\neg (P\lor \neg P)\). Per farlo possiamo pensare di riutilizzarla, in modo da riottenere \(\bot\) usando \(\neg E\), e a quel punto scaricarla con \(\RAA\) per ottenere \(P \lor \neg P\).

\begin{prooftree}
  \AxiomC{\({[\neg P]}_1\)}
  \RuleLabel{\(\lor I^r\)}
  \UnaryInfC{\(P \lor \neg P\)}
  \AxiomC{\({[\neg(P\lor\neg P)]}_2\)}
  \RuleLabel{\(\neg E\)}
  \BinaryInfC{\(\bot\)}
  \RuleLabel{\(\RAA_1\)}
  \UnaryInfC{\(P\)}
  \RuleLabel{\(\lor I^l\)}
  \UnaryInfC{\(P \lor \neg P\)}
  \AxiomC{\({[\neg(P\lor\neg P)]}_2\)}
  \RuleLabel{\(\neg E\)}
  \BinaryInfC{\(\bot\)}
  \RuleLabel{\(\RAA_2\)}
  \UnaryInfC{\(P \lor \neg P\)}
\end{prooftree}

Questo ci mostra che nonostante la regola di deduzione finale sia \(\RAA\) anche partendo da un altra regola di deduzione (\(\lor I\)) siamo riusciti a dimostrare lo stesso \(P \lor \neg P\).

\subsection*{Dimostrazione di \((P \to Q) \leftrightarrow (\neg P \lor Q)\).}

Come abbiamo già visto nell'\Namedcref{ex:notable_tautologies}, \((P \to Q) \leftrightarrow (\neg P \lor Q)\) è una tautologia, quindi ci aspettiamo di poterla dimostrare. Come abbiamo già visto per dimostrare un equivalenza basta dimostrare entrambe le implicazioni che la compongono per poi concludere usando \(\land I\).

Andiamo quindi a dimostrare separatamente \((P \to Q) \to (\neg P \lor Q)\) e \((\neg P \lor Q) \to (P \to Q)\). Banalmente però, sfruttando \(\to I\) possiamo passare al dover dimostrare i sequenti \(P \to Q \vdash \neg P \lor Q\) e \(\neg P \lor Q \vdash P \to Q\), il che ci permette di non doverci preoccupare di scaricare l'assunzione di \(P \to Q\) e \(\neg P \lor Q\) rispettivamente, e quindi di poterle usare liberamente come assunzioni.

\paragraph{Dimostrazione di \(P \to Q \vdash \neg P \lor Q\).}

Dalla dimostrazione precedente, abbiamo visto che un modo per dimostrare formule con \(\lor\) come connettivo principale è quello di usare \(\RAA\).

\begin{prooftree}
  \AxiomC{\(\bot\)}
  \RuleLabel{\(\RAA_1\)}[\({[\neg(\neg P \lor Q)]}_1\)]
  \UnaryInfC{\(\neg P \lor Q\)}
\end{prooftree}

Come già detto numerose volte, per dimostrare \(\bot\) l'idea più promettente è quella di usare \(\neg E\), per la quale ci servono due formule, una la negazione dell'altra. Dall'assunzione di \(P \to Q\) possiamo dedurre \(Q\) assumendo \(P\). Quindi speriamo di poter usare \(\neg E\) con \(Q\) e \(\neg Q\).

\begin{prooftree}
  \AxiomC{\(P\)}
  \AxiomC{\(P \to Q\)}
  \RuleLabel{\(\to E\)}
  \BinaryInfC{\(Q\)}
  \AxiomC{\(\neg Q\)}
  \RuleLabel{\(\neg E\)}
  \BinaryInfC{\(\bot\)}
  \RuleLabel{\(\RAA_1\)}[\({[\neg(\neg P \lor Q)]}_1\)]
  \UnaryInfC{\(\neg P \lor Q\)}
\end{prooftree}

A questo punto dobbiamo scaricare o dimostrare sia \(P\) che \(\neg Q\), poiché \(P \to Q\) non è necessario che venga scaricata.
Dall'assunzione scaricabile di \(\neg(\neg P \lor Q)\) però, dall'equivalenza semantica con \(P \land \neg Q\), ci aspettiamo di poter dedurre sia \(P\) che \(\neg Q\) da \(\land E\).

Iniziamo con \(\neg Q\). Per poter introdurre una negazione, possiamo utilizzare \(\neg I\), che ci permette di dedurre \(\neg Q\) a partire da \(\bot\) scaricando l'assunzione di \(Q\).

\begin{prooftree}
  \AxiomC{\(P\)}
  \AxiomC{\(P \to Q\)}
  \RuleLabel{\(\to E\)}
  \BinaryInfC{\(Q\)}
  \AxiomC{\(\bot\)}
  \RuleLabel{\(\neg I_2\)}[\({[Q]}_2\)]
  \UnaryInfC{\(\neg Q\)}
  \RuleLabel{\(\neg E\)}
  \BinaryInfC{\(\bot\)}
  \RuleLabel{\(\RAA_1\)}[\({[\neg(\neg P \lor Q)]}_1\)]
  \UnaryInfC{\(\neg P \lor Q\)}
\end{prooftree}

A questo punto per dimostrare questo \(\bot\) dobbiamo introdurre una contraddizione. Avendo a disposizione \(Q\) e \(\neg(\neg P \lor Q)\) possiamo pensare di questa seconda insieme a \(\neg P \lor Q\) per dedurre una contraddizione usando \(\neg E\).

\begin{prooftree}
  \AxiomC{\(P\)}
  \AxiomC{\(P \to Q\)}
  \RuleLabel{\(\to E\)}
  \BinaryInfC{\(Q\)}
  \AxiomC{\({[\neg(\neg P \lor Q)]}_1\)}
  \AxiomC{\(\neg P \lor Q\)}
  \RuleLabel{\(\neg E\)}
  \BinaryInfC{\(\bot\)}
  \RuleLabel{\(\neg I_2\)}[\({[Q]}_2\)]
  \UnaryInfC{\(\neg Q\)}
  \RuleLabel{\(\neg E\)}
  \BinaryInfC{\(\bot\)}
  \RuleLabel{\(\RAA_1\)}
  \UnaryInfC{\(\neg P \lor Q\)}
\end{prooftree}

Ma \(\neg P \lor Q\) è deducibile da \(Q\) tramite \(\lor I^r\), permettendoci di chiudere questo ramo di deduzione.

\begin{prooftree}
  \AxiomC{\(P\)}
  \AxiomC{\(P \to Q\)}
  \RuleLabel{\(\to E\)}
  \BinaryInfC{\(Q\)}
  \AxiomC{\({[\neg(\neg P \lor Q)]}_1\)}
  \AxiomC{\({[Q]}_2\)}
  \UnaryInfC{\(\neg P \lor Q\)}
  \RuleLabel{\(\neg E\)}
  \BinaryInfC{\(\bot\)}
  \RuleLabel{\(\neg I_2\)}
  \UnaryInfC{\(\neg Q\)}
  \RuleLabel{\(\neg E\)}
  \BinaryInfC{\(\bot\)}
  \RuleLabel{\(\RAA_1\)}
  \UnaryInfC{\(\neg P \lor Q\)}
\end{prooftree}

Ci rimane quindi solo da dimostrare \(P\). Per farlo però possiamo applicare un ragionamento simile a quello fatto per \(\neg Q\), ma utilizzando \(\RAA\) invece di \(\neg I\) per dedurre \(P\) a partire da \(\bot\) scaricando l'assunzione di \(\neg P\).

\begin{prooftree}
  \AxiomC{\({[\neg(\neg P \lor Q)]}_1\)}
  \AxiomC{\({[\neg P]}_3\)}
  \RuleLabel{\(\lor I^r\)}
  \UnaryInfC{\(\neg P \lor Q\)}
  \RuleLabel{\(\neg E\)}
  \BinaryInfC{\(\bot\)}
  \RuleLabel{\(\RAA_3\)}
  \UnaryInfC{\(P\)}
  \AxiomC{\(P \to Q\)}
  \RuleLabel{\(\to E\)}
  \BinaryInfC{\(Q\)}
  \AxiomC{\({[\neg(\neg P \lor Q)]}_1\)}
  \AxiomC{\({[Q]}_2\)}
  \UnaryInfC{\(\neg P \lor Q\)}
  \RuleLabel{\(\neg E\)}
  \BinaryInfC{\(\bot\)}
  \RuleLabel{\(\neg I_2\)}
  \UnaryInfC{\(\neg Q\)}
  \RuleLabel{\(\neg E\)}
  \BinaryInfC{\(\bot\)}
  \RuleLabel{\(\RAA_1\)}
  \UnaryInfC{\(\neg P \lor Q\)}
\end{prooftree}

L'unica assunzione non scaricata che ci rimane è \(P \to Q\), e quindi abbiamo dimostrato \(P \to Q \vdash \neg P \lor Q\).

\paragraph{Dimostrazione di \(\neg P \lor Q \vdash P \to Q\).}

Ancora una volta, per dimostrare \(P \to Q\) è ragionevole pensare di usare \(\to I\) per ottenere \(P \to Q\) a partire da \(Q\) scaricando l'assunzione di \(P\).

\begin{prooftree}
  \AxiomC{\(Q\)}
  \RuleLabel{\(\to I_1\)}[\({[P]}_1\)]
  \UnaryInfC{\(P \to Q\)}
\end{prooftree}

Ora però dobbiamo dimostrare \(Q\) avendo a disposizione \(\neg P \lor Q\) e \(P\) come assunzioni. Essendoci un \(\lor\) che contiene \(Q\) come operando, è possibile pensare di utilizzare \(\lor E\) per dedurre \(Q\) a partire da \(\neg P \lor Q\), ma per poterlo applicare è necessario dimostrare \(Q\) sia assumendo \(\neg P\) che assumendo \(Q\).

\begin{prooftree}
  \AxiomC{\(\neg P \lor Q\)}
  \AxiomC{\(Q\)}
  \AxiomC{\(Q\)}
  \RuleLabel{\(\lor E_2\)}[\({[\neg P]}_2,\ {[Q]}_2\)]
  \TrinaryInfC{\(Q\)}
  \RuleLabel{\(\to I_1\)}[\({[P]}_1\)]
  \UnaryInfC{\(P \to Q\)}
\end{prooftree}

Per quanto riguarda dimostrare \(Q\) assumendo \(Q\), questo è immediato.
\begin{prooftree}
  \AxiomC{\(\neg P \lor Q\)}
  \AxiomC{\({[Q]}_2\)}
  \UnaryInfC{\(Q\)}
  \AxiomC{\(Q\)}
  \RuleLabel{\(\lor E_2\)}[\({[\neg P]}_2\)]
  \TrinaryInfC{\(Q\)}
  \RuleLabel{\(\to I_1\)}[\({[P]}_1\)]
  \UnaryInfC{\(P \to Q\)}
\end{prooftree}

Ora rimane da dimostrare \(Q\) assumendo \(\neg P\). Per farlo, possiamo utilizzare \(\bot E\), poiché tra le assunzioni scaricabili abbiamo sia \(P\) che \(\neg P\), e quindi possiamo dedurre \(\bot\) usando \(\neg E\), per poi dedurre \(Q\) da \(\bot\).
\begin{prooftree}
  \AxiomC{\(\neg P \lor Q\)}
  \AxiomC{\({[Q]}_2\)}
  \UnaryInfC{\(Q\)}
  \AxiomC{\({[\neg P]}_2\)}
  \AxiomC{\({[P]}_1\)}
  \RuleLabel{\(\neg E\)}
  \BinaryInfC{\(\bot\)}
  \RuleLabel{\(\bot E\)}
  \UnaryInfC{\(Q\)}
  \RuleLabel{\(\lor E_2\)}
  \TrinaryInfC{\(Q\)}
  \RuleLabel{\(\to I_1\)}
  \UnaryInfC{\(P \to Q\)}
\end{prooftree}

In questo esempio, abbiamo potuto osservare come anche le regole di inferza più controintuitive e complesse, come \(\bot E\) e \(\lor E\) possano trovare applicazione.

\section{Proprietà del calcolo deduttivo}

Abbimo visto come il calcolo di deduzione naturale sia un sistema formale che ci permette sintatticamente di ottenere nuove formule a partire da un insieme di premesse, ma per poterlo usare in modo efficace è necessario assicurarsi che esso sia in grado di catturare la conseguenza logica.

In questo senso, andremo a dimostrare varie proprietà del calcolo di deduzione naturale, che lo mettono in relazione con la conseguenza logica, per fare in modo da garantire che questa tecnica di manipolazione sitattica riesca a mantenere valore semantico, e quindi che se una formula è dimostrabile, allora è anche valida, e viceversa.

Iniziamo con una proprietà strettamente relativa al calcolo che ci permette di ampliare arbitrariamente l'insieme di premesse da cui dedurre una formula, senza perdere la possibilità di dedurla.

\begin{theorem}[namedlabel={thm:monotonicity}{Monotonicità}]{Monotonicità di \(\vdash\)}
  Dato un insieme di formule \(\Gamma\) e due formule \(\phi\) e \(\psi\), si ha che
  \[\Gamma \vdash \phi \implies \Gamma \cup \{\psi\} \vdash \phi\]
\end{theorem}

È importante notare che questa proprità risulta solo rafforzata nel caso vengano aggiunte formule che rendono \(\Gamma\) inconsistente, poiché in questo caso sarà possibile dedurre qualsiasi formula a partire da \(\Gamma\) mediante \(\bot E\). Questa è una caratteristica tipica della logica classica, e non è condivisa da tutte le logiche, come ad esempio le logiche paraconsistenti, in cui l'inconsistenza non porta alla trivialità, ma viena in qualche modo contenuta permettendo di parlare del resto dell'insieme di formule anche in presenza di contraddizioni.

\begin{note}{}
  Questa proprietà di generalizza in \(\Gamma \vdash \phi \implies \Gamma \cup \Gamma' \vdash \phi\), applicandola più volte. La dimostrazione di questa versione più generale però segue direttamente da quella della versione più semplice, e quindi non è necessario farla esplicitamente.
\end{note}

\begin{proof}
  Per definizione di deduzione, \(\Gamma \vdash \phi\) se e solo se esiste un albero di derivazione ben formato radice etichettata con \(\phi\) e le cui foglie o sono etichettate con formule di \(\Gamma\) o sono assunzioni scaricabili. 

  In tale definizione non c'è alcun obbligo di usare ogni formula di \(\Gamma\) come etichetta di una foglia. Pertanto, aggiungendo all'insieme delle premesse \(\Gamma\) una nuova formula \(\psi\) non usata nell'albero di derivazione, lo stesso identico albero di derivazione rimane valido per dimostrare \(\Gamma \cup \{\psi\} \vdash \phi\).
\end{proof}

Andiamo quindi a vedere una proprietà analoga a quella vista per la conseguenza logica nel \Namedcref{thm:semantic_deduction}, il vero e proprio \Namedonlyhref{thm:deduction}.

\begin{theorem}[namedlabel={thm:deduction}{Teorema di Deduzione}]{Teorema di Deduzione}
  Dato un insieme di formule \(\Gamma\) e due formule \(\phi\) e \(\psi\), si ha che:
  \[\Gamma, \psi \vdash \phi \iff \Gamma \vdash \psi \to \phi\]
\end{theorem}

\begin{proof}
  Andiamo a dimostrare entrambi i versi.
  
  Iniziamo con \(\Gamma, \psi \vdash \phi \implies \Gamma \vdash \psi \to \phi\).
  Per definizione di deduzione, esiste un albero che dimostra \(\Gamma, \psi \vdash \phi\). Allora è sufficiente applicare la regola \(\to I\) in radice di tale albero, scaricando l'assunzione \(\psi\) e ottenendo così un albero che dimostra \(\Gamma \vdash \psi \to \phi\).

  Nell'altro verso, \(\Gamma \vdash \psi \to \phi \implies \Gamma, \psi \vdash \phi\), l'idea è duale.
  Per definizione di deduzione, esiste un albero che dimostra \(\Gamma \vdash \psi \to \phi\). A questo punto, aggiungendo l'assunzione \(\psi\) all'albero, possiamo applicare la regola \(\to E\) per ottenere un albero che dimostra \(\Gamma, \psi \vdash \phi\).
\end{proof}

Già da questi primi due teoremi è possibile notare un forte parallelismo tra la conseguenza logica e la deduzione, la quale verrà ulteriormente rafforzata dai prossimi teoremi, per poi essere dimostrata in modo definitivo con i teoremi di correttezza e completezza.

\begin{theorem}{Regola del taglio}
  Dato un insieme di formule \(\Gamma\) e due formule \(\phi\) e \(\psi\), si ha che:
  
  Se \(\Gamma, \psi \vdash \phi\) e \(\Gamma' \vdash \psi\) allora \(\Gamma \cup \Gamma' \vdash \phi\)
\end{theorem}

Intuitivamente, questa regola ci dice se abbiamo un albero di derivazione che dimostra un'assunzione, possiamo sostituire a ogni foglia etichettata con quell'assunzione l'albero di derivazione che dimostra quell'assunzione, mantenendo una derivazione valida.


\begin{proof}
  Possiamo formalmente dimostrare questa regola tramite il teorema di deduzione e la monotonicità di \(\vdash\).
  Infatti per monotonicità valgono sia \(\Gamma, \Gamma', \psi \vdash \phi\), che \(\Gamma \cup \Gamma'\vdash \psi\).
  A questo punto, usando il teorema di deduzione, otteniamo \(\Gamma\cup \Gamma' \vdash \psi \to \phi\).

  A questo punto, possiamo costruire un albero di derivazione con \(\phi\) in radice dedotta tramite \(\to E\), i cui due sottoalberi sono l'albero relativo a \(\Gamma \cup \Gamma' \vdash \psi\) e l'albero relativo a \(\Gamma\cup \Gamma' \vdash \psi \to \phi\).
\end{proof}

Prima di spostarci a dimostrare le proprietà di correttezza e completezza, è utile introdurre formalmente l'equivalente sintattico del concetto di insieme insoddisfacile, ovvero quello di insieme inconsistente.

\begin{definition}{Inconsistenza}
  Un insieme di formule \(\Gamma\) è \keyword{inconsistente} se \(\Gamma \vdash \bot\).

  Equivalentemente, \(\Gamma\) è inconsistente, per definizione, se \(\Gamma \vdash \phi\), \(\forall \phi \in \PL\).
\end{definition}

Vedere che le due definizioni sono equivalenti è immediato.
Infatti, se \(\Gamma \vdash \bot\), possiamo applicare \(\bot E\) alla derivazione di \(\bot\), ottenendo così \(\Gamma \vdash \phi\) per ogni formula \(\phi\). Viceversa, se \(\Gamma \vdash \phi\) per ogni formula \(\phi\), allora in particolare si ha che \(\Gamma \vdash \bot\), poiché \(\bot\) è una formula di queste \(\phi\).

Veniamo dunque alle proprietà di correttezza e completezza, che intuitivamente esprimono relazioni di inclusione tra la conseguenza logica e la deduzione. 

\begin{definition}{Completezza}
  Un sistema di deduzione è \keyword{completo} se, per ogni insieme di formule \(\Gamma\) e ogni formula \(\phi\), si ha che:
  \[\Gamma \models \phi \implies \Gamma \vdash \phi\]
\end{definition}

\begin{definition}{Correttezza}
  Un sistema di deduzione è \keyword{corretto} se, per ogni insieme di formule \(\Gamma\) e ogni formula \(\phi\), si ha che:
  \[\Gamma \vdash \phi \implies \Gamma \models \phi\]
\end{definition}

Essendo che sia \(\vdash \subseteq 2^{\PL}\times\PL\) che \(\models \subseteq 2^{\PL}\times\PL\), è possibile pensare a queste due proprietà come a delle relazioni di inclusione tra i due insiemi. In particolare, diremo che il calcolo è corretto se \(\vdash \subseteq \models\), e completo se \(\models \subseteq \vdash\).

Intuitivamente, la completezza di un calcolo ci dice che tutto ciò che è conseguenza logica è anche deducibile, e quindi che il calcolo è in grado di catturare tutta la conseguenza logica. La correttezza invece ci dice che tutto ciò che è deducibile è anche conseguenza logica, e quindi che il calcolo non permette di dedurre formule che non sono conseguenza logica.

Queste due proprietà prese in isolamento però non ci dicono molto, poiché è possibile soddifarle in modo triviale. Per la completezza, ad esempio, è sufficiente prendere come relazione di deduzione l'intero insieme \(2^{\PL}\times\PL\), in modo da poter dedurre qualsiasi formula a partire da qualsiasi insieme di formule, e quindi soddifare la condizione di completezza, ma non quella di correttezza. Per la correttezza invece è sufficiente prendere come relazione di deduzione l'insieme vuoto, in modo da non poter dedurre nessuna formula a partire da nessun insieme di formule, e quindi soddifare la condizione di correttezza, ma non quella di completezza.

Prima di poter dimostrare però che il calcolo di deduzione naturale è corretto e completo, sarà necessario introdurre alcune definizioni e risultati intermedi, che ci permetteranno di dimostrare più facilmente queste due proprietà.

Iniziamo col definire un concetto che ci aiuterà nella dimostrazione di correttezza.

\begin{definition}{Denotazione}
  Data una formula \(\phi\), la \keyword{denotazione} di \(\phi\) è l'insieme di tutti i modelli che soddisfano \(\phi\), ovvero:
  \[\den{\phi} = \set{m \in 2^{\AP}}[ m \models \phi]\]

  Dato un insieme di formule \(\Gamma\), la \keyword{denotazione} di \(\Gamma\) è l'insieme di tutti i modelli che soddisfano tutte le formule di \(\Gamma\), ovvero:
  \[\den{\Gamma} = \set{m \in 2^{\AP}}[ m \models \Gamma] = \bigcap_{\phi \in \Gamma} \den{\phi}\]
\end{definition}

Possiamo dunque procedere alla dimostrazione della correttezza del calcolo di deduzione naturale, che intuitivamente ci dice che tutto ciò che è provabile è semanticamente fondato, e che quindi il calcolo di deduzione naturale ci permette di dedurre solo formule che sono conseguenza logica delle premesse.

\begin{theorem}{Correttezza del calcolo di deduzione naturale}
  Per ogni insieme di formule \(\Gamma\) e ogni formula \(\phi\) si ha che:
  \[\Gamma \vdash \phi \implies \Gamma \models \phi\]
\end{theorem}

\begin{proof}

  Per definzione di \(\Gamma\vdash\phi\), esiste un albero di deduzione con radice \(\phi\) e con foglie che sono o assunzioni scaricate o formule di \(\Gamma\).

  Dunque dimostrare la tesi equivale a dimostrare che, per ogni albero di deduzione in deduzione naturale per \(\phi\) da \(\Gamma\), \(\Gamma\models\phi\).

  Possiamo facilmente vederlo per induzione sulla struttura (o altezza) dell'albero di deduzione.
  
  Il caso base è una deduzione di altezza \(0\), ovvero \(\phi \in \Gamma\). In questo caso, \(\Gamma\models\phi\) è ovviamente verificato.

  A questo punto sia l'altezza maggiore di \(0\). In questo caso avremo che l'albero avrà in radice \(\phi\) dedotta da una certa regola \(\Rules\) e con figli \(\psi_1,\ldots,\psi_k \in \PL\), ognuno di questi a sua volte radice di un sottoalbero di altezza strettamente minore di quella dell'albero di deduzione di \(\phi\). Per ipotesi induttiva, \(\Gamma\models\psi_i\) per ogni \(i=1,\ldots,k\).

  A questo punto quindi ci basta mostrare con un analisi per casi che, per ogni regola di deduzione naturale, se \(\set{\psi_1,\ldots,\psi_k} \vdash_{\Rules} \phi\) allora \(\set{\psi_1,\ldots,\psi_k} \models_{\Rules} \phi\).

  \begin{description}
    \item[\(\to E\).] Abbiamo da ipotesi induttiva che \(\Gamma\models\psi\) e \(\Gamma\models\psi\to\phi\).
      Ovvero, ogni modello che soddisfa \(\Gamma\) soddisfa anche \(\psi\) e \(\psi\to\phi\).
      
      Dalla semantica dell'implicazione, abbiamo che, se un modello \(m\) soddisfa \(\psi\to\phi\) allora o non soddisfa \(\psi\) o, se la soddisfa, allora soddisfa \(\phi\). Ma essendo che \(m\) soddisfa \(\psi\), allora \(m\) deve soddisfare anche \(\phi\). Dunque ogni modello che soddisfa \(\Gamma\) soddisfa anche \(\phi\), e quindi \(\Gamma\models\phi\).
    \item[\(\neg E\).] Abbiamo da ipotesi induttiva che \(\Gamma\models\neg\psi\) e \(\Gamma\models\psi\). Dunque ogni modello che soddisfa \(\Gamma\) soddisfa anche \(\neg\psi\) e \(\psi\). Per semantica del \(\neg\), allora \(m\) non soddisfa \(\psi\), ma abbiamo detto che \(m\) soddisfa \(\psi\), assurdo. Dunque non esiste alcun modello che soddisfa \(\Gamma\), e dunque \(\Gamma\) è insoddisfacibile.

    Una definizione alternativa di conseguenza logica che avevamo dato era che \(\Gamma\models\phi\) se \(\Gamma\cup\set{\neg\phi}\) è insoddisfacibile. Ma allora \(\Gamma\models\bot\) se \(\Gamma\cup\set{\neg \bot}\) è insoddisfacibile. Ma \(\neg\bot\) è tautologicamente vero, dunque \(\Gamma\cup\set{\neg \bot}\) è insoddisfacibile se e solo se \(\Gamma\) è insoddisfacibile. Cosa che abbiamo appena dimostrato nel caso in cui \(\Gamma\models\neg\psi\) e \(\Gamma\models\psi\). Dunque \(\Gamma\models\bot\).

    \item[\(\RAA\).] Abbiamo da ipotesi induttiva che \(\Gamma,\neg\phi\models\bot\). Ovvero \(\Gamma\cup\set{\neg\phi}\) è insoddisfacibile. Dunque \(\Gamma\models\phi\).
    \item[\(\lor E\).] Abbiamo da ipotesi induttiva che \(\Gamma,\psi_1\models\phi\), \(\Gamma,\psi_2\models\phi\) e \(\Gamma\models\psi_1\lor\psi_2\).
    
    Possiamo uniformare le assunzioni usando il teorema di deduzione semantica, ottenendo \(\Gamma\models\psi_1\to\phi\) e \(\Gamma\models\psi_2\to\phi\). Dunque ogni modello che soddisfa \(\Gamma\) soddisfa anche \(\psi_1\to\phi\), \(\psi_2\to\phi\) e \(\psi_1\lor\psi_2\). Ma essendo che \(m\) soddisfa \(\psi_1\lor\psi_2\), allora \(m\) soddisfa almeno una tra \(\psi_1\) e \(\psi_2\). Se \(m\) soddisfa \(\psi_1\), soddisfando \(\psi_1\to\phi\) soddisfa anche \(\phi\). Se \(m\) soddisfa \(\psi_2\), soddisfacendo \(\psi_2\to\phi\) soddisfa anche \(\phi\). Dunque \(m\) soddisfa \(\phi\) in ogni caso.
    
    Alternativamente, possiamo dimostrare questo punto usando \(\den{\phi}\) e \(\den{\Gamma}\).

    Essendo che \(\Gamma\models \psi_1 \lor \psi_2\), ogni modello che soddisfa \(\Gamma\) soddisfa anche \(\psi_1\lor\psi_2\). Dunque \(\forall m \in \den{\Gamma}, m \in \den{\psi_1\lor\psi_2}\). Alternativamente, \(\den{\Gamma}\subseteq \den{\psi_1\lor\psi_2}\). Con lo stesso ragionamento, dalla semantica della disgiunzione, \(\forall m \in \den{\psi_1\lor\psi_2}, m \in \den{\psi_1}\cup \den{\psi_2}\). Inoltre, abbiamo che \(\den{\Gamma}\subseteq\den{\psi_1\to\phi}\) e \(\den{\Gamma}\subseteq\den{\psi_2\to\phi}\), e dalla semantica dell'implicazione, si ha che \(\den{\Gamma} \cap \den{\psi_1}\subseteq\den{\phi}\) e \(\den{\Gamma}\cap\den{\psi_2}\subseteq\den{\phi}\), poiché ogni modello che soddisfa \(\Gamma\) e soddisfa \(\psi_1\) o \(\psi_2\) deve soddisfare anche \(\phi\). In particolare quindi
    \[(\den{\Gamma}\cap\den{\psi_1})\cup(\den{\Gamma}\cap\den{\psi_2})\subseteq \den{\phi}\]
    dalla distributività dell'intersezione rispetto all'unione, si ha che
    \[\den{\Gamma} \cap (\den{\psi_1}\cup\den{\psi_2})\subseteq\den{\phi}\]
    Avendo già mostrato però che \(\den{\Gamma} \subseteq \den{\psi_1}\cup\den{\psi_2}\) la loro intersezione è proprio \(\den{\Gamma}\),
    dunque \(\den{\Gamma}\subseteq \den{\phi}\), ovvero ogni modello che soddisfa \(\Gamma\) soddisfa anche \(\phi\), e dunque \(\Gamma\models\phi\).
    \item[] Per questioni di brevità non verranno trattate tutte le regole, ma solo quelle più interessanti. Gli altri casi sono analoghi a quelli già trattati, e possono essere dimostrati con un ragionamento simile a quello fatto per le regole già trattate.
  \end{description}
\end{proof}

Conclusa la dimostrazione di correttezza, possiamo passare al verso più interessante, ovvero quello di completezza, che ci dice che tutto ciò che è conseguenza logica è anche deducibile, e quindi che il calcolo di deduzione naturale è in grado di catturare tutta la conseguenza logica.

Questo risulta un po' più complesso da dimostrare, poiché, a differenza della correttezza che chiedeva semplicemente di verificare che le prove generate dal calcolo esibissero una certa proprietà, ovvero l'essere conseguenze logiche delle premesse, la completezza ci chiede di verificare che l'intero spazio delle conseguenze logiche sia coperto dal calcolo, e quindi non sarà più sufficiente verificare in maniera strutturale le prove.

Per dimostrare il teorema di correttezza infatti faremo uso di un paio di risultati intermedi, che ci permetteranno di dimostrare la completezza in modo più semplice.

\begin{lemma}[namedlabel={lem:lindenbaum}{Lindenbaum}]{Lindenbaum}
  Sia \(\Sigma\subseteq\PL\) un insieme consistente di formule. Allora è possibile costruire un insieme di formule \(\Sigma'\) tale che \(\Sigma\subseteq\Sigma'\) e \(\Sigma'\) è massimalmente consistente, ovvero per ogni formula \(\psi\), o \(\psi\in\Sigma'\) o \(\neg\psi\in\Sigma'\).
\end{lemma}

\begin{proof}
  Consideriamo un'arbitraria enumerazione delle formule di \(\PL\), ovvero una sequenza \(\phi_1,\phi_2,\ldots\) che contiene tutte le formule di \(\PL\).
  
  Si costruisca dunque una sequenza di insiemi di formule \(\Sigma_0,\Sigma_1,\ldots\) tale che \(\Sigma_0 = \Sigma\) e che
  
  \[ \Sigma_{i+1} = \begin{cases}
                      \Sigma_i\cup\set{\phi_{i+1}} & \text{se } \Sigma_i\cup\set{\phi_{i+1}} \text{ è consistente} \\
                      \Sigma_i & \text{altrimenti}
                    \end{cases}
                    \]

  In questo modo è garantito che, se \(\Sigma\) è consistente, per induzione ogni \(\Sigma_i\) è consistente.

  Definiamo allora \(\Sigma^* = \bigcup_{i\geq 0}\Sigma_i\). Dobbiamo allora mostrare che \(\Sigma^*\) è massimalmente consistente.

  È importante notare che, nonostante \(\Sigma^*\) sia il limite di una successione di cui ogni elemento presenta una proprietà, non è detto che \(\Sigma^*\) presenti quella proprietà, e quindi è necessario dimostrarlo esplicitamente. Si pensi ad esempio a una qualsiasi successione numerica diventante. Ogni elemento di questa successione ha la proprietà di essere finito, ma il limite di questa successione è infinito, e quindi non presenta la proprietà di essere finito.

  Per dimostrare dunque che \(\Sigma^*\) è massimalmente consistente dovremo prima verificare che \(\Sigma^*\) è consistente, e poi verificare che \(\Sigma^*\) è massimale.

  Per assurdo, assumiamo che non sia consistente, ovvero \(\Sigma^*\vdash\bot\). Per definizione di deduzione, deve esistere una deduzione naturale di \(\bot\) da \(\Sigma^*\) che chiamiamo \(\pi\). Una deduzione però è un oggetto finito, dunque \(\pi\) contiene solo un numero finito di foglie, e dunque solo un numero finito di assunzioni non scaricate. Dunque dev'esistere un \(\Gamma\) finito tale che \(\Gamma\subset\Sigma^*\) e \(\Gamma\vdash\bot\). Ma essendo che ogni \(\Sigma_i\) contiene tutti i \(\Sigma_j\) con \(j<i\), posso prendere un \(i\) pari al massimo indice di una formula di \(\Gamma\) rispetto all'enumerazione scelta. Dunque \(\Gamma\subseteq\Sigma_i\), e dunque \(\Sigma_i\vdash\bot\), ovvero \(\Sigma_i\) non è consistente, assurdo.\footnote{Ciò che in questo caso ci permette di dimostrare che il limite mantiene la proprietà degli elementi della successione è il fatto che la deduzione è un oggetto finito, e quindi un'eventuale violazione della proprietà di consistenza da parte del limite deve essersi verificata già in un punto della successione, e quindi in un elemento della successione.}

  Passando dunque alla massimalità, dobbiamo dimostrare che, per ogni formula \(\psi\), o \(\psi\in\Sigma^*\) o \(\neg\psi\in\Sigma^*\).
  Possiamo però notare che, se esiste \(\psi\not\in\Sigma^*\) tale che \(\Sigma^*\cup\set{\psi}\) è consistente, allora \(\Sigma^*\) non è massimalmente consistente, poiché \(\neg\psi\not\in\Sigma^*\), altrimenti si potrebbe dedurre \(\bot\) da \(\Sigma^*\cup\set{\psi}\) usando \(\neg E\), e quindi \(\Sigma^*\cup\set{\psi}\) sarebbe inconsistente, assurdo. Di conseguenza per mostrare che \(\Sigma^*\) è massimalmente consistente è sufficiente mostrare che, per ogni formula \(\psi\) tale che \(\psi\not\in\Sigma^*\), \(\Sigma^*\cup\set{\psi}\) è inconsistente.\footnote{Ad essere precisi, qui vale l'equivalenza tra le condizioni. Infatti, se \(\Sigma^*\) è massimalmente consistente, allora per ogni formula \(\psi\not\in\Sigma^*\) si ha che necessariamente \(\neg\psi\in\Sigma^*\), e quindi \(\Sigma^*\cup\set{\psi}\) è inconsistente.}

  Dall'enumerazione delle formule di \(\PL\) scelta in precedenza, deve esistere un \(i \in \N\) tale che \(\phi_i = \psi\).
  Per costruzione di \(\Sigma^*\), se \(\psi\not\in\Sigma^*\) allora \(\Sigma_{i-1}\cup\set{\psi}\) è inconsistente, altrimenti \(\psi = \phi_i\) sarebbe stata aggiunta a \(\Sigma_{i-1}\) e quindi sarebbe stata aggiunta anche a \(\Sigma^*\).

  Ma allora \(\Sigma_{i-1}\cup\set{\psi}\vdash\bot\). Essendo \(\Sigma_{i-1}\cup\set{\psi} \subseteq \Sigma^*\cup\set{\psi}\), dal \Namedcref{thm:monotonicity} si ha che \(\Sigma^*\cup\set{\psi}\vdash\bot\), assurdo.
\end{proof}

A questo punto, possiamo a dimostrare una proprietà più generale, che ci permetterà di dimostrare la completezza del calcolo di deduzione naturale in modo più semplice, ovvero che un insieme di formule consistente è soddisfacibile.

Per fare ciò pero, dobbiamo andare ad introdurre un concetto simile a quello di denotazione, ma in un certo senso duale, ovvero la \keyword{teoria} di un modello.

\begin{definition}{Teoria}
  Dato un modello \(m\), la \keyword{teoria} di \(m\) è l'insieme di tutte le formule che sono soddisfatte da \(m\), ovvero:
  \[\Th{m} = \set{\phi \in \PL }[ m \models \phi]\]
\end{definition}

In un certo senso, sia la denotazione che la teoria collegano il mondo sintattico a quello semantico tramite una sorta di proprietà di chiusura. Infatti, la denotazione di una formula è l'insieme di tutti i modelli che la soddisfano, e quindi è un insieme di modelli, mentre la teoria di un modello è l'insieme di tutte le formule che sono soddisfatte da quel modello, e quindi è un insieme di formule.

Andiamo dunque a vedere qualche proprietà delle teorie di modelli, che ci permetteranno di costruire un modello a partire da un insieme di formule consistente.

\begin{proposition}{}
  Sia \(m\) un modello e \(\Th{m}\) la sua teoria. Allora \(\Th{m}\) è massimalmente consistente.
\end{proposition}

\begin{proof}
  Per definizione di teoria, \(\Th{m} = \set{\phi\in\PL}[m\models\phi]\). Dunque, se \(\Th{m}\) fosse inconsistente, allora \(\Th{m}\vdash\bot\), e quindi dalla correttezza del calcolo di deduzione naturale, \(\Th{m}\models\bot\), e quindi \(m\models\bot\), assurdo. Dunque \(\Th{m}\) è consistente.

  Per dimostrare che \(\Th{m}\) è massimale, dobbiamo dimostrare che, per ogni formula \(\psi\), o \(\psi\in\Th{m}\) o \(\neg\psi\in\Th{m}\). Ma essendo \(m\) un modello, per ogni formula \(\psi\), \(m\) soddisfa esattamente una tra \(\psi\) e \(\neg\psi\). Dunque, se \(m\models\psi\), allora \(m\not\models\neg\psi\), e quindi \(\neg\psi\notin\Th{m}\), ma \(m\models\psi\), dunque \(\psi\in\Th{m}\). Se invece \(m\models\neg\psi\), allora \(m\not\models\psi\), e quindi \(\psi\notin\Th{m}\), ma \(m\models\neg\psi\), dunque \(\neg\psi \in \Th{m}\).
\end{proof}

\begin{proposition}[label=prop:max-consistent-set-are-theory]{}
  Sia \(\Sigma^* \subseteq \PL\) massimalmente consistente. Allora esiste un modello \(m\) tale che \(\Sigma^* = \Th{m}\).
\end{proposition}

\begin{proof}
  Dato \(\Sigma^*\) massimalmente consistente, dobbiamo costruire un modello \(m\) tale che \(\Sigma^* = \Th{m}\).
  Per fare ciò, basta definire \(m = \Sigma^*\cap\AP\), poiché essendo \(\Sigma^*\) massimalmente consistente, contiene tutte e sole le proposizioni atomiche che non portano a inconsistenza. Infatti le uniche proposizioni atomiche che non possono essere contenute in \(\Sigma^*\) sono quelle che, se aggiunte a \(\Sigma^*\), rendono \(\Sigma^*\) inconsistente, ovvero appaiono in qualche formula che è già presente in \(\Sigma^*\) e che le richiede con assegnamento falso per essere soffisfatte.

  A questo punto, dobbiamo mostrare che \(\Sigma^* = \Th{m}\), ovvero che ogni formula \(\phi\), \(\phi \in \Sigma^* \iff m \models \phi\).

   Per fare ciò, possiamo procedere per induzione strutturale sulle formule.

  Il caso base è una formula atomica \(p\). Se \(p\in\Sigma^*\), allora \(p\in m\) per costruzione, e quindi \(m\models p\) e viceversa.
  A questo punto, sia \(\phi\) una formul di altezza maggiore di \(0\). In questo caso, \(\phi\) è costruita a partire da formule di altezza minore tramite l'applicazione di una regola di formazione. L'ipotesi induttiva è dunque che per ogni \(\psi\) di altezza minore di \(\phi\),\(\psi \in \Sigma^*\) se e solo se \(m\models\psi\).
  \begin{description}
    \item[\(\phi = \neg\psi_1\)] Essendo \(\Sigma^*\) massimalmente consistente, se \(\neg\psi_1\in\Sigma^*\) allora \(\psi_1\notin\Sigma^*\). Per ipotesi induttiva, \(m\) non soddisfa \(\psi_1\), dunque \(m\) soddisfa \(\neg\psi_1\), ovvero \(m\models\phi\). Viceversa, se \(m\models\neg\psi_1\) allora \(m\) non soddisfa \(\psi_1\), dunque, per ipotesi induttiva, \(\psi_1\notin\Sigma^*\). Essendo \(\Sigma^*\) massimalmente consistente, se \(\psi_1\notin\Sigma^*\) allora \(\neg\psi_1\in\Sigma^*\), dunque \(\phi = \neg\psi_1\in\Sigma^*\).
    \item[\(\phi = \psi_1\to\psi_2\)] Essendo \(\Sigma^*\) massimalmente consistente, se \(\psi_1\to\psi_2\in\Sigma^*\) allora \(\psi_1\notin\Sigma^*\) o \(\psi_2\in\Sigma^*\). Se \(\psi_1\notin\Sigma^*\), allora \(m\) non soddisfa \(\psi_1\), dunque \(m\) soddisfa \(\psi_1\to\psi_2\), ovvero \(m\models\phi\). Se invece \(\psi_2\in\Sigma^*\), allora \(m\) soddisfa \(\psi_2\), dunque \(m\) soddisfa \(\psi_1\to\psi_2\), ovvero \(m\models\phi\). Viceversa, se \(m\models\psi_1\to\psi_2\) allora o \(m\) non soddisfa \(\psi_1\) o \(m\) soddisfa \(\psi_2\). Se \(m\) non soddisfa \(\psi_1\), allora per ipotesi induttiva \(\psi_1\notin\Sigma^*\), dunque, essendo \(\Sigma^*\) massimalmente consistente, \(\neg\psi_1\in\Sigma^*\). Ma essendo \(\Sigma^*\) massimalmente consistente, deve contenere anche \(\psi_1\to\psi_2\), ovvero \(\phi\), poiché \(\neg\psi_1\vdash\psi_1\to\psi_2\) e \(\Sigma^*\) è chiuso per deduzione. Analogamente, se \(m\) soddisfa \(\psi_2\), allora per ipotesi induttiva \(\psi_2\in\Sigma^*\), dunque, essendo \(\Sigma^*\) massimalmente consistente, \(\psi_1\to\psi_2\in\Sigma^*\), ovvero \(\phi\in\Sigma^*\) poiché \(\psi_2\vdash\psi_1\to\psi_2\) e \(\Sigma^*\) è chiuso per deduzione.
    \item[\(\phi = \psi_1\lor\psi_2\)] Essendo \(\Sigma^*\) massimalmente consistente, se \(\psi_1\lor\psi_2\in\Sigma^*\) allora almeno una tra \(\psi_1\) e \(\psi_2\) è presente in \(\Sigma^*\). Se \(\psi_1\in\Sigma^*\), allora \(m\) soddisfa \(\psi_1\), dunque \(m\) soddisfa \(\psi_1\lor\psi_2\), ovvero \(m\models\phi\). Se invece \(\psi_2\in\Sigma^*\), allora \(m\) soddisfa \(\psi_2\), dunque \(m\) soddisfa \(\psi_1\lor\psi_2\), ovvero \(m\models\phi\). Viceversa, se \(m\models\psi_1\lor\psi_2\) allora \(m\) soddisfa almeno una tra \(\psi_1\) e \(\psi_2\). Se \(m\) soddisfa \(\psi_1\), allora per ipotesi induttiva \(\psi_1\in\Sigma^*\), dunque, essendo \(\Sigma^*\) massimalmente consistente, \(\psi_1\lor\psi_2\in\Sigma^*\), ovvero \(\phi\in\Sigma^*\). Se invece \(m\) soddisfa \(\psi_2\), allora per ipotesi induttiva \(\psi_2\in\Sigma^*\), dunque, essendo \(\Sigma^*\) massimalmente consistente, \(\psi_1\lor\psi_2\in\Sigma^*\), ovvero \(\phi\in\Sigma^*\).
    \item[\(\phi = \psi_1\land\psi_2\)] Essendo \(\Sigma^*\) massimalmente consistente, se \(\psi_1\land\psi_2\in\Sigma^*\) allora sia \(\psi_1\) che \(\psi_2\) sono presenti in \(\Sigma^*\). Dunque \(m\) soddisfa sia \(\psi_1\) che \(\psi_2\), dunque \(m\) soddisfa \(\psi_1\land\psi_2\), ovvero \(m\models\phi\). Viceversa, se \(m\models\psi_1\land\psi_2\) allora \(m\) soddisfa sia \(\psi_1\) che \(\psi_2\). Dunque per ipotesi induttiva sia \(\psi_1\) che \(\psi_2\) sono presenti in \(\Sigma^*\), dunque, essendo \(\Sigma^*\) massimalmente consistente, \(\psi_1\land\psi_2\in\Sigma^*\), ovvero \(\phi\in\Sigma^*\).
  \end{description}
\end{proof}

\begin{lemma}[namedlabel={lem:model-existence}{Esistenza del modello}]{Esistenza del modello}
  Dato \(\Sigma\subseteq\PL\), si ha che:
  
  \(\Sigma\) è consistente \(\implies\Sigma\) è soddisfacibile.
\end{lemma}

\begin{proof}
  Per definizione di consistenza \(\Sigma \not\vdash \bot\), ovvero non esiste alcun albero di deduzione con radice \(\bot\) e con foglie che sono o assunzioni scaricate o appartengono a \(\Sigma\). Essere soddisfacibile significa che esiste un modello \(m\) tale che \(m\models\Sigma\).

  Dunque dobbiamo andare a costruire un modello \(m\) che soddisfa \(\Sigma\) sapendo che le formule di \(\Sigma\) sono consistenti.

  Ma a questo punto, se \(\Sigma\) è consistente, allora per il \Namedcref{lem:lindenbaum} è possibile costruire un insieme di formule \(\Sigma^*\) tale che \(\Sigma\subseteq\Sigma^*\) e \(\Sigma^*\) è massimalmente consistente. Dunque, per la \Cref{prop:max-consistent-set-are-theory}, esiste un modello \(m\) tale che \(\Sigma^* = \Th{m}\). Dunque, essendo \(\Sigma\subseteq\Sigma^*\), allora \(\Sigma\subseteq\Th{m}\), e quindi \(m\models\Sigma\), ovvero \(m\) soddisfa tutte le formule di \(\Sigma\), e quindi \(\Sigma\) è soddisfacibile.
\end{proof}

A questo punto, siamo finalmente in grado di dimostrare la completezza del calcolo di deduzione naturale, che ci dice che tutto ciò che è conseguenza logica è anche deducibile, e quindi che il calcolo di deduzione naturale è in grado di catturare tutta la conseguenza logica.

\begin{theorem}{Completezza del calcolo di deduzione naturale}
  Per ogni insieme di formule \(\Gamma\) e ogni formula \(\phi\) si ha che:
  \[\Gamma \models \phi \implies \Gamma \vdash \phi\]
\end{theorem}

\begin{proof}
  Come già visto precedentemente, abbiamo che \(\Gamma\models\phi \iff \Gamma \cup \set{\neg\phi}\) è insoddisfacibile.

  Inoltre, possiamo facilmente vedere che \(\Gamma\vdash \phi \iff \Gamma \cup \set{\neg\phi}\) è inconsistente.

  Infatti, se \(\Gamma \cup \set{\neg\phi}\) è inconsistente allora \(\Gamma \cup \set{\neg\phi}\vdash \bot\), ovvero esiste un'albero di deduzione per \(\bot\) da \(\Gamma \cup \set{\neg\phi}\). Dunque, se scarichiamo l'assunzione \(\neg\phi\) con \(\RAA\), otteniamo un albero di deduzione per \(\phi\) da \(\Gamma\), ovvero \(\Gamma\vdash\phi\).

  Viceversa, se \(\Gamma\vdash\phi\) allora esiste un albero di deduzione per \(\phi\) da \(\Gamma\). Se aggiungiamo a questo albero di deduzione un'assunzione \(\neg\phi\) e applichiamo \(\neg E\), allora otteniamo un albero di deduzione per \(\bot\) da \(\Gamma \cup \set{\neg\phi}\), ovvero \(\Gamma \cup \set{\neg\phi}\vdash\bot\), e quindi \(\Gamma \cup \set{\neg\phi}\) è inconsistente.

  Di conseguenza, \(\Gamma\models\phi \implies \Gamma\vdash\phi\) è equivalente a \(\Gamma \cup \set{\neg\phi}\) è insoddisfacibile \(\implies\) \(\Gamma \cup \set{\neg\phi}\) è inconsistente.

  A questo punto, possiamo dimostrare la proprietà di completezza per contrapposizione, ovvero che, preso un qualsiasi insieme \(\Gamma\), \(\Gamma\) consistente \(\implies\) \(\Gamma\)soddisfacibile.

  Ma per fare ciò è sufficiente il \Namedonlyhref{lem:model-existence}, che ci dice che un insieme di formule consistente è soddisfacibile. Di conseguenza, se \(\Gamma\cup\set{\neg\phi}\) è insoddisfacibile, allora è inconsistente, ovvero se \(\Gamma\models\phi\) allora \(\Gamma\vdash\phi\).
\end{proof}